\section{Experiments}
\label{sec:experiments}

\subsection{Experimental Setup}

\paragraph{Datasets.} We evaluate on four Amazon product review domains with varying characteristics:

\begin{table}[h]
\centering
\small
\caption{Dataset statistics.}
\label{tab:datasets}
\begin{tabular}{lrrrl}
\toprule
Domain & Test Users & Candidates & Avg Title Len & Characteristics \\
\midrule
Beauty & 973 & 101 & 37 chars & Short titles, niche categories \\
Books & 10,000 & 101 & 49 chars & Rich titles, diverse genres \\
Electronics & 10,000 & 101 & 15 chars & Technical specs, short titles \\
Movies & 10,000 & 101 & 21 chars & Very short titles, genre-driven \\
\bottomrule
\end{tabular}
\end{table}

Each test instance consists of 1 positive item (ground-truth next interaction) and 100 randomly sampled negative items, following the standard same-candidate evaluation protocol~\cite{krichene2020sampled}.

\paragraph{Backbone Model.} All methods use Qwen3-8B as the unified language model backbone. This ensures fair comparison: differences in performance reflect the \emph{scoring methodology}, not the model capacity.

\paragraph{Baselines.} We compare against 8 official LLM-based recommendation methods, all adapted to use Qwen3-8B:

\begin{itemize}
    \item \textbf{ELMRec}~\cite{elmrec}: Graph-enhanced LLM recommendation
    \item \textbf{IRLLRec}~\cite{irllrec}: Intent-aware retrieval with LLM
    \item \textbf{LLM2Rec}~\cite{llm2rec}: LLM-to-recommendation with SASRec
    \item \textbf{LLMEmb}~\cite{llmemb}: LLM embedding for recommendation
    \item \textbf{LLM-ESR}~\cite{llmesr}: LLM-enhanced sequential recommendation
    \item \textbf{ProEx}~\cite{proex}: Profile-based explanation recommendation
    \item \textbf{ProMax}~\cite{promax}: Profile-maximized recommendation
    \item \textbf{RLMRec}~\cite{rlmrec}: Reinforcement learning meets LLM recommendation
\end{itemize}

All baselines use their official code or official-code-level implementations with default hyperparameters, adapted only for the unified backbone and data interface.

\paragraph{Metrics.} We report HR@$k$ (Hit Rate), NDCG@$k$ (Normalized Discounted Cumulative Gain) for $k \in \{5, 10, 20\}$, and MRR (Mean Reciprocal Rank).

\paragraph{Implementation Details.} Our method uses vLLM for batch inference with \texttt{max\_model\_len=2048}, \texttt{temperature=0.1}, \texttt{max\_new\_tokens=100}, and prefix caching enabled. The prompt uses the 5 most recent history items and up to 200 characters of candidate description.

\subsection{Main Results}

\begin{table*}[t]
\centering
\small
\caption{Main same-candidate results. Each domain has 10,000 users and 101
candidates per user. We show \method{} and the strongest official baseline by
NDCG@10; the full evidence ledger contains all eight official baselines per
domain. Bold marks the better value between the two shown rows.}
\label{tab:main_results}
\begin{tabular}{llrrrrrrr}
\toprule
Domain & Method & HR@5 & HR@10 & HR@20 & NDCG@5 & NDCG@10 & NDCG@20 & MRR \\
\midrule
\multirow{2}{*}{Sports}
& Best official (LLMEmb) & 0.2124 & 0.3384 & 0.4900 & 0.1389 & 0.1795 & 0.2177 & 0.1539 \\
& \textbf{\method} & \textbf{0.2745} & \textbf{0.3819} & \textbf{0.5172} & \textbf{0.1985} & \textbf{0.2329} & \textbf{0.2670} & \textbf{0.2076} \\
\midrule
\multirow{2}{*}{Toys}
& Best official (LLMEmb) & 0.2499 & 0.3505 & 0.4866 & 0.1725 & 0.2049 & 0.2391 & 0.1814 \\
& \textbf{\method} & \textbf{0.3172} & \textbf{0.3964} & \textbf{0.5059} & \textbf{0.2452} & \textbf{0.2708} & \textbf{0.2983} & \textbf{0.2503} \\
\midrule
\multirow{2}{*}{Home}
& Best official (LLMEmb) & 0.1079 & 0.1856 & 0.3169 & 0.0690 & 0.0939 & 0.1267 & 0.0901 \\
& \textbf{\method} & \textbf{0.1561} & \textbf{0.2264} & \textbf{0.3505} & \textbf{0.1098} & \textbf{0.1324} & \textbf{0.1635} & \textbf{0.1259} \\
\midrule
\multirow{2}{*}{Tools}
& Best official (LLMEmb) & 0.1365 & 0.2257 & 0.3637 & 0.0875 & 0.1159 & 0.1505 & 0.1065 \\
& \textbf{\method} & \textbf{0.1937} & \textbf{0.2696} & \textbf{0.3931} & \textbf{0.1419} & \textbf{0.1661} & \textbf{0.1970} & \textbf{0.1559} \\
\bottomrule
\end{tabular}
\end{table*}


\subsection{Calibration Analysis}

To validate our motivating observation, we compute calibration metrics for all methods that produce numerical scores (4 out of 8 baselines output raw scores; the remaining 4 produce only ranked lists without scores).

\begin{table}[t]
\centering
\small
\caption{Calibration metrics for methods with raw score outputs. All use Qwen3-8B backbone, 101-candidate protocol. Lower ECE/MCE/Brier is better; higher AUROC is better.}
\label{tab:calibration}
\begin{tabular}{ll|cccc}
\toprule
Domain & Method & AUROC$\uparrow$ & ECE$\downarrow$ & MCE$\downarrow$ & Brier$\downarrow$ \\
\midrule
\multirow{5}{*}{Beauty}
& \textbf{Ours} & 0.613 & \textbf{0.237} & \textbf{0.728} & \textbf{0.098} \\
& IRLLRec & \textbf{0.622} & 0.479 & 0.926 & 0.279 \\
& LLMEmb & 0.624 & 0.508 & 0.924 & 0.311 \\
& RLMRec & 0.601 & 0.518 & 0.924 & 0.319 \\
& LLM2Rec & 0.496 & 0.494 & 0.947 & 0.294 \\
\midrule
\multirow{5}{*}{Books}
& LLMEmb & \textbf{0.778} & 0.493 & 0.882 & 0.291 \\
& LLM2Rec & 0.701 & 0.623 & 0.902 & 0.434 \\
& IRLLRec & 0.615 & 0.446 & 0.922 & 0.246 \\
& RLMRec & 0.590 & 0.406 & \textbf{0.934} & \textbf{0.214} \\
\midrule
\multirow{4}{*}{Electronics}
& LLMEmb & \textbf{0.658} & 0.506 & 0.931 & 0.306 \\
& IRLLRec & 0.563 & 0.444 & 0.945 & 0.244 \\
& RLMRec & 0.545 & 0.388 & 0.954 & 0.199 \\
& LLM2Rec & 0.543 & 0.315 & 0.968 & 0.140 \\
\midrule
\multirow{4}{*}{Movies}
& LLMEmb & \textbf{0.722} & 0.546 & \textbf{0.915} & 0.355 \\
& LLM2Rec & 0.665 & 0.662 & 0.917 & 0.498 \\
& IRLLRec & 0.618 & 0.394 & 0.929 & 0.222 \\
& RLMRec & 0.608 & 0.386 & 0.937 & 0.206 \\
\bottomrule
\end{tabular}
\vspace{-2mm}
\end{table}


Key findings:
\begin{enumerate}
    \item All baselines exhibit severe miscalibration (ECE $> 0.31$) across all domains.
    \item Maximum Calibration Error exceeds 0.88 universally, indicating that certain confidence regions are completely unreliable.
    \item Our method achieves the lowest ECE in all tested domains, confirming that task-grounded probability scoring produces better-calibrated outputs.
\end{enumerate}
